\documentclass[twoside]{article}

\usepackage{ustj}

\addbibresource{mss.bib}

\newcommand{\authorname}{K. Afonin}
\newcommand{\authorpatp}{\patp{dozreg-toplud}}
\newcommand{\affiliation}{Urbit Foundation}

%  Make first page footer:
\fancypagestyle{firststyle}{%
\fancyhf{}% Clear header/footer
\fancyhead{}
\fancyfoot[L]{{\footnotesize
              %% We toggle between these:
              Manuscript submitted for review.\\
              % {\it Urbit Systems Technical Journal} I:1 (2024):  1–3. \\
              ~ \\
              Address author correspondence to \authorpatp.
              }}
}
%  Arrange subsequent pages:
\fancyhf{}
\fancyhead[LE]{{\urbitfont Urbit Systems Technical Journal}}
\fancyhead[RO]{Subject Knowledge Analysis}
\fancyfoot[LE,RO]{\thepage}

%%MANUSCRIPT
\title{Subject Knowledge Analysis}
\author{\authorname~\authorpatp \\ \affiliation}
\date{}

\begin{document}

\maketitle
\thispagestyle{firststyle}

\begin{abstract}
\end{abstract}

% We will adjust page numbering in final editing.
\pagenumbering{arabic}
\setcounter{page}{1}

\tableofcontents

\section{Introduction}

To date, the Nock combinator calculus or instruction set (depending on one's viewpoint) has been evaluated using either a tree-walking interpreter or a bytecode interpreter.  The former has been found straightforward to implement but regrettably slow in practice.  Following a suggestion by \patp{ritpub-sipsyl} to incorporate better information about what is known about the subject into the compilation process, Urbit engineers elaborated \textsc{ska}, or Subject Knowledge Analysis.  This static analysis pass is performed on a subject-formula pair, producing a call graph and a subject mask.  The latter can be used to determine if two subject-formula pairs are equivalent, allowing the reuse of analysis results.  The call graph can be used to introduce direct calls in the bytecode interpreter, as well as to perform compile-time jet matching.  \textsc{ska} unlocks ample gains in practical Nock performance, as well as enabling future optimizations and analyses.

\section{Dynamism of Nock}

The Nock 2 operator is one of the two opcodes that are used for Nock evaluation.  From the Nock specification:

\begin{lstlisting}[style=listingcode]
*[a 2 b c]          *[*[a b] *[a c]]
\end{lstlisting}

In plain words, we evaluate two operands of Nock 2 against the given subject, then evaluate Nock with the product of the first operand as a subject, and the product of the second operand as a formula.

The compound Nock 9 operator is defined in terms of other Nock operators:

\begin{lstlisting}[style=listingcode]
*[a 9 b c]          *[*[a c] 2 [0 1] 0 b]
\end{lstlisting}

Evaluation of the formula from the second operator yields a noun, from which a formula to be evaluated is pulled with Nock 0.  That formula is then evaluated against that noun, which is often referred to as a ``core''.

Since Nock 9 is essentially a macro for Nock 2, we can focus our attention on Nock 2.  This operator is equivalent to \lstinline[style=inlinecode]{eval} in other languages:  execution of
dynamically generated code.  Unlike other languages, Nock does not have a notion of a static call, and \emph{all} function calls are implemented in terms of eval.  In practice, however, we can discover formula products \lstinline[style=inlinecode]{*[a c]} for almost all Nock 2 sites, given the subject.

Examples when we cannot do that, such as receiving code
to evaluate over the wire, or from the outer virtualization context via Nock 12, or as a product of a complex expression where we cannot reliably execute it at compile time, are rare.  In practice, these exceptions are limited to either evaluating
dynamically-supplied code or running the Hoon compiler at run time, e.g. in vase mode.\footnote{See \patp{lagrev-nocfep}, \emph{Urbit Systems Technical Journal} vol.\ 2 issue 1, pp.~131–153, for more information on vase mode.}  For all other cases, e.g. Arvo, Gall agents, the subject's code contents do not often change and the result of static analysis is likely to be reused until a kernel upgrade or Gall agent upgrade respectively.

Thus ``Subject Knowledge Analysis'' (\textsc{ska}) is the name for this static analysis procedure:  given a pair of a subject and a formula, we construct the call graph for all Nock 2 sites for which we have enough information, yielding call graph information and subject mask.  A new subject-formula pair where the formula is equal to the analyzed formula, and the parts of the subject included in the mask are equal to the same parts from the previously analyzed subject, is considered to be equivalent.  Consider analyzing the subject-formula pair:

\begin{lstlisting}[style=listingcode]
[[dec 1] [%9 2 %10 [6 %0 3] %0 2]]
\end{lstlisting}

Then a pair \lstinline[style=listingcode]{[[dec 2] [%9 2 %10 [6 %0 3] %0 2]]} is equivalent to the first one in terms of the call graph information: the different subnouns of the subject are never used as code by the \lstinline[style=listingcode]{$} arm of the \lstinline[style=listingcode]{dec} core. A pair of a minimized subject and a formula thus constructs an object equivalent to a function in other languages, with the masked out parts of the subject being the function's arguments.  Further in the text, we will refer to pairs (minimized subject)-formula as functions that we discovered during the analysis.

\section{Motivation}

In prior work on \textsc{ska} by \patp{ritpub-sipsyl} and \patp{master-morzod}, the analysis was considered as a first step in compiling Nock to a static single assignment intermediary representation language (\textsc{ssa} IR) to introduce direct calls in Nock. In their absence, Nock registerization, or any kind of optimization/analysis for that matter is limited by the lack of knowledge about the call target and the call product: in the \textsc{ssa} IR example, we would not know which parts of a call's
subject are actually used by the callee program.

For the current work, a simpler goal was set: the introduction of direct calls
and compile-time jet matching in the Vere bytecode interpreter.  Currently, the
interpreter has to, first, look up a Nock bytecode program in the bytecode
cache, and second, it has to dynamically perform a jet matching routine at each
Nock 9 call site (although there has been work done to ameliorate the situation,
see https://github.com/frodwith/nockets).  With \textsc{ska}, the bytecode programs could
have instructions for direct calls into other bytecode programs, and it could
also check if a given Nock 2 site matches some jet information it accumulated,
and adding a jet driver directly into the bytecode stream if the match was
obtained.

## Core algorithm

To analyze a subject-formula pair, we run what is essentially a partial Nock
interpreter to propagate known information to formulas deeper in the formula
tree, and to accumulate call graph information.  Instead of a noun, the subject
is a partial noun described with the type \lstinline[style=listingcode]{$sock}:

\begin{lstlisting}[style=listingcode]
+$  cape  $~(| $@(? [cape cape]))
+$  sock  $~([| ~] [=cape data=*])
\end{lstlisting}

That is, `cape` is a mask that describes the shape of the known parts of a given
noun, and `sock` is a pair of a mask and data, where unknown parts are stubbed
with `0` or `~`.

The treatment of the partial noun by that interpreter is mostly trivial:

- Autocons conses the partial products, normalizing the mask and data if
 necessary;
- Nock 0 grabs a subnoun from \lstinline[style=listingcode]{$sock} if it can, otherwise it returns unknown
 result `[| ~]`;
- Nock 1 returns a fully known result;
- Actual computations (Nock 3, 4, 5) return unknown result `[| ~]` -- we do not
 attempt to run all code at compile time, we just want to evaluate code-
 generating expressions;
- Nock 6 produces an intersection of the results of branches, that is, only
 parts of the products that are known and equal between branches;
- Nock 7 composes formulas: the product of the first formula is passed as a
 partial subject to the second one;
- Nock 10 edits the product of one formula with the product of the other.  If
 the recipient noun is not known, and the donor noun is known at least
 partially, cells will be created that lead to the edited part.
- Nock 8, 9 are desugared in terms of Nock 7, Nock 2, and autocons
- Nock 11 have their dynamic hint formula products dropped, \lstinline[style=listingcode]{%spot} hints are
 used for debugging/verbose printouts, \lstinline[style=listingcode]{%fast} hints are used to accumulate
 cold state;
- Nock 12 of the virtual Nock produces an unknown result.

In the case of Nock 2, the operands are evaluated, and then we check if the
formula at this site is fully known.  If it is not known, the call is \emph{indirect} -- we cannot make any assumptions about its product, so we return
`[| ~]`.  If it is known, we do two things:

1. We record information on subject usage as code for each call below us in the
 stack, including the root call.  This step is needed to produce a minimized
 subject for each call, which is necessary to avoid redoing the work for the
 functions we already analyzed.
2. We enter a new frame of the partial interpreter, executing/analyzing new
 sock-formula pair we just got, returning the result of that analysis as a
 product.

To track code usage of function calls' subjects, we pair the partial noun with
a structure that describes the flow of data from the functions below us in the
stack to the new callsite.  That data structure will be described later in
the performance section.

## Loop handling

The most complicated part of the algorithm is the handling of loops.  Unlike the
regular interpreter, where only one branch is executed at a time, the partial
interpreter enters both, so a lack of loop handling would quickly lead to an
infinite cycle in the analysis, be it a simple decrement loop or mutual
recursion in the Hoon compiler.

### Loop detection

One way to detect loops would be to simply compare sock-formula pairs that we
have on the stack with the one that we are about to evaluate.  That, however,
would not work with functions that, for example, recursively construct a list,
as the subject sock would change from iteration to iteration.  So instead of
a simple equality check, we need to see if the new subject (further referred to
as \emph{kid} subject, or subject of the \emph{kid call}) nests under the subject of the
supposedly matching function call on the stack (\emph{parent} subject, or subject of
the \emph{parent call}), given the current information we have about the usage of
parent subject as code.

The problem is that the code usage information for the parent is not complete
until we analyzed all its callees/descendants.  Therefore, that loop guess would
have to be validated when we return from the loop, and before that, no analysis
of function calls in that loop can be finalized.  We can see that we have three
kinds of function calls: ones that are not a part of any cycle, and can be
finalized immediately, ones that are in the middle of a cycle and cannot be
finalized unless the entry point into that cycle is finalized, and finally,
those that are the entry points into a cycle, whose finalization also means
finalizing all members of a cycle.  Another interesting aspect is that we do not
know which kind of a function call we have on entry, we know it only on return.
For the non-finalizable calls, we also do not know which call finalization they
depend on until that loop entry is finalized, since we could always discover a
new loop call that points deeper into the stack.

### Call graph topology

First, let us describe in more detail what the abstract call graph is and how it
is traversed/discovered in \textsc{ska}.

The call graph is a directed graph with the vertices being the functions and the
edges being function calls.  The graph has a root, which is the top-level
function that is being analyzed.  We include function calls from callers to
callees, even if the calls are conditional: from the point of view of the call
graph topology, conditional calls of functions B and C performed by a function A
are indistinguishable from consecutive calls, i.e. these two Hoon expressions
would yield the same call graph shape (but the other information could, of
course, be different):

```hoon
?:  condition
  (func-a x)
(func-b x)
::
(func-b (func-a x))
```

We can see that, when we execute Nock, we traverse the call graph in
depth-first, head-first order, except we skip some edges if they are
conditional.  As we enter other functions, the path from the root vertex to the
current one forms the computational stack.

The abstract call graph may contain cycles.  Figure 1 shows the abstract call
graph for Nock obtained from compiling `(dec 42)` against `..ride` subject.
Firstly `dec` expression is evaluated, which pulls `+dec` arm from a core in the
subject.  With the product of that expression, which is the `dec` core, pinned
to the subject, we edit that core with the argument `42`, then we pull the arm
\lstinline[style=listingcode]{$} from that core.  The \lstinline[style=listingcode]{$} arm call either returns the decremented value or
calls itself, incrementing the counter.

| ![](dec_abstract.png) | 
|:--:| 
| *Figure 1.  Abstract call graph for `(dec 42)`* |

When we traverse that call graph, we cannot simply follow the edge from \lstinline[style=listingcode]{$:dec}
to itself, as doing so would cause an infinite loop in the algorithm.  Instead,
when we detect a loop call, or a \emph{backedge}, we immediately return an unknown
result from the child call.  As we traverse the call graph, we enumerate the
vertices, as shown in Figure 2.

| ![](dec-ska.png) | 
|:--:| 
| Figure 2.  Call graph for `(dec 42)` in the middle of \textsc{ska}.  Vertices 0x2 and 0x3 form a loop (strongly connected component, SCC) with vertex 0x2 being the entry point (magenta), and vertex 0x3 being the latest (i.e. rightmost, deepest) vertex that belongs to that SCC, also called a \emph{latch} of this SCC (orange).  The dashed gray line denotes the loop assumptions that are to be validated upon exit from the SCC entry point.  Vertex 0x3 is not analyzed through.|

Figure 2 also demonstrates that the \textsc{ska} representation of the call graph is
acyclic: back-pointing edges are replaced with a forward-pointing edge to a
new vertex that we do not analyze through.  Together with the fact that both
branches are traversed in Nock 6 analysis, \textsc{ska} call graph traversal becomes a
genuine depth-first, head-first traversal.

### Loops as strongly connected components

From the point of view of the abstract, cyclical call graph, a set of functions
that cannot be finalized unless all functions from this set are finalized forms
a strongly-connected component (SCC) of that call graph.  Thus, handling of
loops in \textsc{ska} becomes a question of handling SCCs.

What follows is the algorithm that the author developed to detect and manage
SCCs during the analysis.  A reader well-versed in graph theory and/or PL theory
will likely recognise it as some well-known algorithm, as it is not that
special.  The author's ignorance of that subject, however, prevents him from
adding a reference to that algorithm in the literature.

As the call graph is traversed, we keep track of the stack of mutually
independent SCCs (that is, SCCs whose unions are not SCCs) by keeping track of
the entry point (the earliest vertex that has a backedge pointing to it) and the
\emph{latch} (the latest vertex that has a backedge originating from it).

Suppose that we encounter a new backedge.  To determine whether it forms a new
SCC, or belongs to an already existing one, it is sufficient to compare the 
parent vertex enumeration label with the latch of the latest SCC on the loop
stack.  If the parent vertex is greater than the latch, a new SCC data structure 
is pushed on the stack; otherwise, the top SCC's entry and latch are updated, and
then that SCC is repeatedly merged with the one beneath it if the connection
condition is satisfied: the entry of the updated SCC is greater than or equal to
the latch of the preceding SCC.

Let us prove that this comparison is sufficient.  Recall that the enumeration
labels are assigned in depth-first, head-first order.  If the backedge parent is
equal to the latch of the SCC, the kid is a part of the SCC since the kid is
reachable by an SCC member, in this case, the parent, and a member of SCC is
reachable from kid via the backedge (figure 3).

| ![](parent-eq-latch.png) | 
|:--:| 
| *Figure 3.  New backedge points to the latch of previous SCC 0x1, extending that SCC* |

If the parent label is greater than the latch, then either they are on the same
root path or not.  In both cases, no member of the SCC can be reached from the
parent, making the parent, kid, and all vertices in between a new SCC.
These scenarios are illustrated in Figure 4.

| ![](parent-gth-latch.png) | 
|:--:| 
| *Figure 4.  Illustration of a new backedge pointing to a vertex with a label greater than the previous SCC's latch label, creating a new SCC with 0x2 as entry and 0x4 as latch on the left, and 0x3 as an entry and 0x4 as a latch on the right.* |

Finally, if the parent label is less than the latch, then either they are on the
same root path or not.

If they are, the latch is the parent's descendant, making the
latest SCC reachable from the parent.  Since the parent and the kid are
necessarily of the same root path, which makes the kid reachable from the latch,
making both the parent and the kid vertices reachable from the SCC, thus making
the kid, parent, and all vertices between them are part of the previous SCC.

If they are not, then the latch still has to be a parent's descendant.  Let us
sketch the proof by disproving the contrary: suppose that the opposite is true.
Then the parent vertex is to the left of the latch, where "A is to the left of
B" means that A is entered before B and A is not B's predecessor.  The latch is
either on the same root path as the kid or it is to the left of the kid vertex,
since the latter is the latest vertex we discovered.  Since "leftness" is
transitive, that means that the parent is to the left of the kid, which is
impossible since they are on the same root path.  Contradiction: the latch and
the parent are on the same root path, making the kid, parent vertices, and
vertices between them parts of the existing SCC.  Illustrations for both
scenarios are shown in Figure 5.

| ![](parent-lth-latch.png) | 
|:--:| 
| *Figure 5.  Illustration of a new backedge pointing to a vertex with a label less than the previous SCC's latch label, extending that SCC and moving the latch to a new position.* |

The condition for SCC merging in case of SCC extension is the same.  In fact SCC
extension can be represented in terms of SCC creation, together with merging
with the topmost SCC in the loop stack.

### Cycle validation, code usage distribution, fix-point search

Since loop validation is deferred until returning from the loop entry point,
subject usage by the loop calls is not recorded until that point.  At loop
validation, we want to distribute subject code usage information of the parent
over the kid's subject provenance, finally making an update we deferred so far.
But by doing so, since the kid's subject may contain nouns from the parent call,
given that the kid is always a parent's descendant in the call graph, we might
update usage information of the parent's subject.  We would have to do that in a
loop until the parent's subject usage converged to the fixpoint value for that
transformation.

For each backedge that we are validating, we perform the following steps:

1. We calculate the parent subject, masked down with the subject code usage
 information;
2. We distribute the cape of that masked subject using kid provenance
 information;
3. We recalculate the new masked-down parent subject.
4. We compare their capes. If they are equal, the fixpoint is found; otherwise,
 return to step 2
5. Check masked-down parent subject and kid subject for nesting.  If they nest,
 the backedge is valid; otherwise, the assumption was wrong, and we will redo
 the analysis of the loop with this parent-kid pair added to a blacklist.

Steps 1 and 3 are necessary because otherwise, if we distribute just the subject
code usage information cape without using it to mask actual data, the cape could
start growing infinitely deep, causing an infinite loop in the fixpoint search.
Step 4 works because the cape distribution along the given provenance is
idempotent.

## Performance: load-bearing design choices and optimizations

### Vastness of nouns

It is important to have a sense of scale when dealing with \textsc{ska}.  Consider, for
example, Urbit's standard library core `zuse`.  How many nouns are there in
total in that library?  In Arvo 409K:

```hoon
> =/  n=*  ..zuse
  |-  ^-  @
  ?@  n  1                    ::  atom is one noun
  ~+                          ::  %memo hint
  +((add $(n -.n) $(n +.n)))  ::  cell is 1 noun + nouns in head + nouns in tail
1.656.521.503.863.284.811.637
```

$1.66 * 10^{21}$ nouns!  "Billions of billions" would be an underestimate.

But Urbit's runtime, Vere, handles this mass of nouns just fine.  We can
estimate that it takes only ~1.2 MB to store that library.  What is going on?
And how could we calculate the number of nouns in a reasonable time with
something that looks like naive tree traversal?

The answer is structural sharing.  The library carries a lot of duplicated
nouns, which are represented as pointers to reference-counted data structures.
The reason we could count the number of nouns with something that seems like
a simple tree traversal is the use of \lstinline[style=listingcode]{%memo} hint that enables memoization of
Nock results, using the subject-formula pair as a key.  Even if the nouns are
large, their comparison is still fast precisely because they are duplicated:
noun equality check is short-circuited on pointer equality.  So even though
pointer equality is not exposed in Nock directly, having nouns be
pointer-equivalent is crucial for performance.  Thus, when developing algorithms
with performance in mind, we need to:

1. Make sure not to disrupt structural sharing whenever possible;
2. Leverage structural sharing by either short-circuiting various
 checks/comparisons with Nock 5 equality test, or by using Nock memoization.

### Provenance tree

We represent the provenance information of the subject with \lstinline[style=listingcode]{$source}:

\begin{lstlisting}[style=listingcode]
+$  source  (lest (lest *))
\end{lstlisting}

where `lest` is a non-empty list.  

Raw nouns in the inner list describe the provenance from a function call's
subject to a use site of that noun.  The value can be:

- 0, if the noun does not come from the subject, e.g., it was created
 by evaluating Nock 1 from the formula;
- a non-zero atom, in which case the noun comes from the given axis of the
 subject;
- a cell, in which case the noun in question is also a cell, and its head and
 tail have provenances that are described by the head and tail of the
 provenance noun respectively.

Nock 1, 3, 4, 5, 12 erases provenance information, Nock 0 gets a sub-provenance
tree in an obvious way, Nock 7 composes provenance calculation, Nock 10 performs
an edit, all similar to the \lstinline[style=listingcode]{$sock} treatment.  What differs significantly is
the handling of the branching operator Nock 6.  What we care about is where the
data \emph{could} come from, so we need to calculate a union of provenances, masked
down to the intersection of produced data.  Instead of creating a tree that
includes all union information, in the spirit of saving structural sharing, we
simply make a list of simple provenances.  Thus `(lest *)` is a union of
provenances in a given call.

The outer list keeps entries per call frame.  Whenever we enter a new call with
Nock 2, we push the new provenance union list `~[1]` onto the provenance stack.
When we return from a call, the stack is popped, and the popped provenance union
list is composed with the second-to-top.  Whenever we need to distribute subject
usage, e.g., when we encounter a direct call, the provenance stack is folded
with composition.

### Provenance tree operations with heuristics

When dealing with provenance information \lstinline[style=listingcode]{$source}, some heuristics were added
to improve performance.  When applying binary functions like composition to two
union lists of provenances, the function would be applied to all pairs, and a
new union list would be constructed, omitting duplicates or other provenances
that nest under some other provenance that is already present in the assembled
list.  The compatibility check function only checks provenance nouns up to ten
cells deep, assuming that they are incompatible beyond that point.  This
heuristic prevents from walking provenance nouns exhaustively and taking too
much time if they are too deep at the expense of possibly having larger union
lists with duplicates, which does not affect the semantics of the algorithm.
In case of Nock 10 provenance edit, the check was disabled altogether if the
lengths of the union lists of donor and recipient nouns multiplied exceeded 100.

Since the provenance union is expressed as a list, there are two ways of
implementing the consing of provenances.  We can either cons every member of the
list of the head provenance with every member of the list of the tail
provenance, yielding a union list with the length equal to the product of the
lengths of the operand lists.  Or we can cons every member of the list of the
head provenance with empty provenance `0`, cons `0` with every member of the list
of the tail provenance, then concatenate the lists.  That way, we end up with a
union list whose length is a sum of the operand lists' lengths.  The algorithm
for consing chooses between two approaches depending on the length of the
operand lists, using the multiplicative approach almost always, since typically
the union list length is one, and switching to the additive "cons via union" if
the union lists are too long.

### Load-bearing comparison checks and \lstinline[style=listingcode]{%memo} hints

In \lstinline[style=listingcode]{$cape} and \lstinline[style=listingcode]{$sock} arithmetic, equality check shortcuts for cape union,
which is heavily used for code usage information updates, and sock intersection,
which is used for Nock 6 product calculation, are load-bearing.  For \lstinline[style=listingcode]{%memo} hints, we use them to short-circuit
Nock memoization hints, these are load-bearing for provenance noun composition
and construction of subject capture mask used for in-algorithm caches described
below.  They also appear to add marginal performance improvements in other
places of the algorithm in case of backtracking (reanalyzing a cycle over and
over due to making incorrect guesses).

### Algorithm caches for finalized and non-finalized calls

To prevent doing the same work over and over, we want to cache the intermediary
analysis results.  For finalized calls, the cached data includes the formula and
the minimized subject, where the latter includes potential subject capture in
addition to code usage.  If a Nock 2 formula gets a cache hit, it updates the
code usage information using the mask from the cache, and returns the product,
which it constructs by taking the cached product sock and composing the cached
union list with current subject provenance, effectively applying the provenance
transformation of the cached call to the new subject.

Caching finalized calls turned out to be not enough, which was demonstrated by
\patp{ritpub-sipsyl} and \patp{master-morzod} at
[Lake Summit](https://www.youtube.com/watch?v=z24zWdHfVfI).  To prevent
reanalyzing cycles in highly mutually-recursive code like the Hoon compiler,
caching of non-finalized calls was also necessary.  Similarly to finalized-call
caching, a potential cache match would check the subject it has for nesting with
the cached subject sock, masked down at cache check time with the currently
available code usage information.  Similarly to loop calls, the code usage
information distribution along the new subject provenance would be deferred to
the cycle validation.  Another point of similarity is the necessity of checking
cycles for overlap and merging in case they do, since a cache hit call would
have a loop call to a member of an SCC to which the cached function belongs.

Unlike loop calls, the fixpoint search is not necessary since the cached call
and the cache hit are never on the same root path, so distributing the code
usage information along the hit call subject provenance would never have an
effect on cached call.  Validation is still necessary, and reanalysis of cycles
due to incorrect non-finalized call cache guesses is currently the biggest
overhead in the analysis.

## Vere integration

Integration of \textsc{ska} analysis into Vere was done in a simplest imaginable way to
produce a demo.  Road struct was updated to include bytecode caches for
\textsc{ska}-produced code and to house the analysis core.  A static hint \lstinline[style=listingcode]{%ska} was
added as an entry point into the analysis-compilation lifecycle.

### Nomm (Nock--)

In addition to producing call graph information, the \textsc{ska} algorithm also produced
desugared, annotated Nock-like code called Nomm (Nock--).  The practical
difference, apart from Nock 9 and Nock 8 expanded into other Nock operators, is
the annotation data in Nock 2.  Nomm version that is handled by the current
version of Vere fork with \textsc{ska}, \lstinline[style=listingcode]{$nomm-1}, is mostly indistinguishable from
\lstinline[style=listingcode]{$nock} proper:

\begin{lstlisting}[style=listingcode]
+$  nomm-1
  $^  [nomm-1 nomm-1]
  $%  [%1 p=*]
      [%2 p=nomm-1 q=nomm-1 info=(unit [less=sock fol=*])]
      [%3 p=nomm-1]
      [%4 p=nomm-1]
      [%5 p=nomm-1 q=nomm-1]
      [%6 p=nomm-1 q=nomm-1 r=nomm-1]
      [%7 p=nomm-1 q=nomm-1]
      [%10 p=[p=@ q=nomm-1] q=nomm-1]
      [%11 p=$@(@ [p=@ q=nomm-1]) q=nomm-1]
      [%12 p=nomm-1 q=nomm-1]
      [%0 p=@]
  ==
\end{lstlisting}

In the \lstinline[style=listingcode]{%2} case, `info=~` denotes an indirect call, and a non-empty case
denotes a direct call with a known formula and masked subject.  Here, the pair
of a minimized subject and a formula is used as a unique identifier for a Nomm
function.

### Road data structures

Two HAMT tables were added to the \lstinline[style=listingcode]{u3a_road} struct: \lstinline[style=listingcode]{dar_p}, a map from
[sock formula] pair to loom offset of \lstinline[style=listingcode]{u3n_prog} bytecode \lstinline[style=listingcode]{u3p(u3n_prog)}, and
\lstinline[style=listingcode]{lar_p}, a map from formula to a list of pairs \lstinline[style=listingcode]{[sock u3p(u3n_prog)]}.  The
former is used during Nomm compilation to bytecode, the latter is used for
searching for bytecode with direct calls, given a formula and a subject.

The Nock bytecode struct was extended with an array of direct call information
structs:

\begin{lstlisting}[style=listingcode_c]
/* u3n_dire: direct call information
  */
typedef struct {
  u3p(u3n_prog)  pog_p;
  u3j_harm*      ham_u;
  c3_l           axe_l;
} u3n_dire;
\end{lstlisting}

where

\begin{itemize}
  \item \lstinline[style=listingcode]{pog_p} is a bytecode program loom offset for the direct call,
  \item \lstinline[style=listingcode]{ham_u} is a nullable pointer to a jet driver if the jet match occurred at compile time,
  \item \lstinline[style=listingcode]{axe_l} is the arm axis for jet driver debugging checks.
\end{itemize}

Another addition to the road struct was a field for the analysis core `+ka`.
A jammed cell of four pre-parsed Hoon ASTs for \textsc{ska} source files was added to
Vere source code, which was unpacked and built into a core when the \textsc{ska} entry
point was executed with `+ka` core not being initialized.

### Analysis lifecycle

Let us follow the execution of a [subject formula] pair using \textsc{ska}.  It starts
with annotating a computation with a \lstinline[style=listingcode]{%ska} hint:

\begin{lstlisting}[style=listingcode]
~>  %ska
\end{lstlisting}

#### \lstinline[style=listingcode]{%ska} hint

When the regular bytecode compiler in Vere encounters a \lstinline[style=listingcode]{%ska} hint, it defers
the compilation of the hinted formula and emits \textsc{ska}-entry point opcode.  When
that opcode is executed, a lookup is performed in \lstinline[style=listingcode]{lar_p}, returning \lstinline[style=listingcode]{u3n_prog*}
on success.  If the lookup failed, the pair of subject and formula is passed to
\lstinline[style=listingcode]{+rout} gate of the \lstinline[style=listingcode]{+ka} core, which returned \lstinline[style=listingcode]{+ka} core with global code
tables and cold state updated.  After that, the code maps were extracted from
the global state, and the \textsc{ska} compiler function was called with the subject,
formula, and the maps as arguments.  The maps include:

\begin{itemize}
  \item \lstinline[style=listingcode]{cole}: a map from \lstinline[style=listingcode]{[sock *]} to \lstinline[style=listingcode]{[path axis]} for cold state accumulation,
  \item \lstinline[style=listingcode]{code}: a map from \lstinline[style=listingcode]{[sock *]} to loom offset of \lstinline[style=listingcode]{u3n_prog} for finalized call caching,
  \item \lstinline[style=listingcode]{fols}: a map from formula to a list of pairs \lstinline[style=listingcode]{[sock u3p(u3n_prog)]} for non-finalized call caching.
\end{itemize}

#### Nomm compilation and execution

The main difference between compiling regular Nock to Vere bytecode and
compiling Nomm is that compiling a Nomm formula with direct calls also requires
compiling all its callees recursively.  A rewrite step was added to handle
cycles properly: on the first pass \lstinline[style=listingcode]{u3_noun} values for \lstinline[style=listingcode]{[sock *]} pairs were
saved in \lstinline[style=listingcode]{u3n_dire} structs, and on the second pass, these values were replaced
with corresponding \lstinline[style=listingcode]{u3p(u3n_prog)} values.

Once the bytecode program for the entry point function and all its descendants
were compiled, the Nock bytecode interpreter would execute this program as any
other program, with the only difference of having direct call opcodes that skip
bytecode cache lookup and jet matching.

### Benchmarks, notable results

The effect of the \lstinline[style=listingcode]{%ska} hint on performance was compared for the Ackermann
function and for a simple algorithm to iterate over a list and compute on its
values.  In both cases, the performance gain was around 1.7x.  The main overhead
in computations with direct calls appears to come from
refcounting/allocating/freeing operations and interpreter overhead (figures 6,
7, download to zoom/search).

| ![alt text](perf_409_debug_999_noska.svg) | 
|:--:| 
| *Figure 6.  Flame graph for Ackermann function computation without \lstinline[style=listingcode]{%ska} hint* |

| ![alt text](perf_409_debug_999.svg) | 
|:--:| 
| *Figure 7.  Flame graph for Ackermann function computation with \lstinline[style=listingcode]{%ska} hint* |

\sloppy
When it comes to the performance of the \textsc{ska} algorithm itself, the results were sometimes surprising.  Paradoxically, analyzing \lstinline[style=listingcode]{+mint:ut} call itself took \lstinline[style=listingcode]{~m1}, while analyzing the entirety of \lstinline[style=listingcode]{hoon.hoon}, including \lstinline[style=listingcode]{+mint:ut}, takes \lstinline[style=listingcode]{~s13}.  The reason appears to be a lack of knowledge in the solo \lstinline[style=listingcode]{+mint:ut} call case, causing excessive backtracking due to wrong non-finalized call cache hits (thus validating \textsc{ska}'s fundamental thesis). In addition, every time the loop analysis was discarded, all caches were also discarded, leading to making more work multiple times when compared to the \lstinline[style=listingcode]{hoon.hoon} case.  In the latter case, no wrong cache hits were made.

\section{Conclusion and future work}

There appear to be two viable future directions for this project.  One, more short-term oriented is proper integration with Vere, e.g., making compiled \textsc{ska} bytecode caches persistent and integrating cold states of Vere and \textsc{ska}.  A longer-term direction would continue \patp{ritpub-sipsyl}'s original vision by implementing \textsc{ssa} IR compilation, which would ameliorate allocation overhead caused by cell churning during temporary core construction for calls and deconstruction to get function arguments.\footnote{Compare the rationale for the implementation of the \lstinline[style=listingcode]{#} hax edit instruction in Nock 4K (\patp{lagrev-nocfep} and \patp{sorreg-namtyv}\emph{Urbit Systems Technical Journal} vol.\ 2 issue 1, pp.~155–190).}\tombstone

\selectlanguage{USenglish}
\printbibliography
\end{document}
